\documentclass[11pt]{article}

\usepackage[margin=1in]{geometry}
\usepackage{amsmath,amssymb,amsthm,mathtools}
\usepackage{enumitem}
\usepackage{booktabs}
\usepackage{tabularx}
\usepackage{microtype}
\usepackage[hidelinks]{hyperref}
\usepackage{xurl}

\newtheorem{definition}{Definition}
\newtheorem{lemma}{Lemma}
\newtheorem{remark}{Remark}

\title{Challenge Coverage Witness Slices for Successor Maintenance:\\[0.5ex]
\large Audience-Minimal Route-and-Branch Certification Deltas for One Coverage Claim}
\author{}
\date{Unified archive note v0.04 (2026-03-21; archive v484)}

\begin{document}
\maketitle

\begin{abstract}
A challenge coverage witness pack can bind one certified coverage entry to the aggregate route-and-branch bundle that certifies that claim as a whole.
A witness core can bind that same entry to the audience-invariant terminal-witness and surface-hook keys reused by every certified audience.
One narrower question remains: what is the minimal \emph{route-and-branch delta} for one audience's use of that already-certified claim?
This note isolates that audience-specific layer.
It defines a compact \emph{challenge coverage witness slice} that binds one coverage entry and one audience to exactly one route key together with the route-specific branch keys needed for that audience alone, while importing common terminal and surface proving keys from the companion witness core.
\end{abstract}

\paragraph{Non-redundancy note.}
Synthesis~39 owns only audience coverage and support-only status, Synthesis~40 owns exact piece-to-segment realization, Synthesis~41 owns exact piece-hook terminal ownership, Synthesis~42 owns the aggregate route-and-branch proving bundle for one coverage claim, Synthesis~44 owns the audience-invariant proving core, Synthesis~54 owns the answer-side terminal citation normal form, Synthesis~74 owns the paired terminal citation discipline, Synthesis~75 owns the paired cutover rule, and Synthesis~17 owns the concrete worked instantiation.
This note owns only the audience-minimal route-and-branch slice of that proving bundle once later-paper stop posture is already fixed. Later notes should reopen it only when the exact audience-specific proving delta is itself live.
\paragraph{Scope.}
This is not a new coverage certificate, route map, branch map, terminal-witness ledger, surface-hook ledger, witness core, or aggregate witness pack.
It is a narrow publication-interface note about recording the smallest route/branch dependency set for one \emph{audience-specific} use of one already-declared coverage claim.

\paragraph{Imports.}
Sentence-family coverage is imported from Synthesis~39.\cite{series:challengecoverage}
Exact surface realization is imported from Synthesis~40.\cite{series:challengesurfacehooks}
Exact terminal ownership is imported from Synthesis~41.\cite{series:challengeterminalwitnesses}
Aggregate coverage witness packs are imported from Synthesis~42.\cite{series:challengecoveragewitnesspacks}
Audience-invariant witness cores are imported from Synthesis~44.\cite{series:challengecoveragewitnesscores}
The answer-side terminal citation normal form and its paired reuse/cutover discipline are imported from Synthesis~54 and Synthesis~74--75 as the default later-paper stop rule around this audience-specific delta owner.\cite{series:challengeanswerreviewmenus,series:pairedterminalcitationdiscipline,series:pairedterminalcutoverrules}
The concrete worked instantiation is imported from Synthesis~17.\cite{series:workedexample}
\paragraph{Exports.}
This note exports (i) challenge coverage witness slices, (ii) an audience-minimality rule, (iii) a route-specific branch rule, (iv) a pack-union rule relating slices back to one aggregate witness pack and one shared witness core, and (v) a local-reopen rule saying that later notes may cite one maintained witness slice only when the exact audience-specific route-and-branch delta is itself live.
Otherwise later terse papers should default to the answer-side terminal citation normal form in Synthesis~54 together with the paired terminal discipline and paired cutover rule in Synthesis~74--75. The audience-invariant proving core remains outside scope here and belongs to Synthesis~44.\cite{series:challengecoveragewitnesscores,series:challengeanswerreviewmenus,series:pairedterminalcitationdiscipline,series:pairedterminalcutoverrules}
\section{Why audience-minimal slices deserve their own note}
An aggregate witness pack is the right object for certifying one coverage claim as a whole.
A witness core is the right object for storing the audience-invariant terminal/surface dependencies reused by every audience.
But later papers often reuse that claim for only one audience at a time: a terse release note need not carry a referee route, and a referee footnote need not carry an auditor route.
If the archive offers only aggregate packs, maintainers either over-cite more route/branch keys than one audience needs, or implicitly prune them ad hoc at reuse time.
So the archive benefits from one small object whose only job is to say which route/branch keys are minimally sufficient for one audience-specific use of one already-certified coverage claim.

\begin{definition}[Challenge coverage witness slice]
Fix one concrete base$\to$successor comparison, one already-declared coverage entry $e$, and one audience $a$ certified for that entry by a companion coverage certificate.
Let $\Pi(e)$ be the companion aggregate witness pack for $e$ and let $\Gamma(e)$ be the companion witness core.
A \emph{challenge coverage witness slice} for $(e,a)$ is a tuple
\[
\Sigma_a(e)=(e,a,f,p,m,\pi,\gamma,r,B)
\]
containing the coverage-entry key $e$, audience $a$, sentence family $f$, semantic piece key $p$, coverage mode $m\in\{\textsf{surface},\textsf{support-only}\}$, the witness-pack key $\pi$, the witness-core key $\gamma$, one route key $r$, and a finite nonempty set $B$ of branch keys witnessed by that route-plus-piece pair.
\end{definition}

\begin{definition}[Audience-minimality rule]
A witness slice satisfies the \emph{audience-minimality rule} if it names exactly one route key $r$, that route belongs to audience $a$, and no route for any other audience appears in the slice.
\end{definition}

\begin{definition}[Route-specific branch rule]
A witness slice satisfies the \emph{route-specific branch rule} if every listed branch in $B$ is witnessed by the same route-plus-piece pair induced by $r$ and $p$, and the slice points to the same witness-core key $\gamma$ already named by the companion aggregate witness pack $\pi$.
\end{definition}

\begin{definition}[Pack-union rule]
A family of audience slices satisfies the \emph{pack-union rule} if the union of their route keys and branch keys over all certified audiences for $e$ recovers exactly the route set and branch set named by the companion aggregate witness pack $\Pi(e)$ up to order, while every such slice points back to the same witness core $\Gamma(e)$.
\end{definition}

\begin{lemma}[Audience slices prevent over-citation drift]
Assume every audience-specific use of a coverage entry is paired with a witness slice satisfying the audience-minimality rule, the route-specific branch rule, and the pack-union rule.
Then later notes may cite only the route-and-branch keys needed for one audience without silently changing either the aggregate dependency bundle that certifies the coverage claim as a whole or the common proving core reused across audiences.
\end{lemma}

\begin{proof}
The audience-minimality rule prevents one audience-specific citation from silently dragging in routes meant only for other audiences.
The route-specific branch rule prevents a slice from drifting away from the exact route-plus-piece walk already certified for that audience.
The pack-union rule keeps all slices anchored to one already-declared aggregate witness pack and one shared witness core, so per-audience minimization cannot drift away from the full certified dependency bundle.
Therefore the archive can support audience-minimal citation practice without weakening machine-checkable certification.\qedhere
\end{proof}

\begin{table}[t]
\centering
\small
\begin{tabularx}{\linewidth}{@{}l l X@{}}
\toprule
Coverage claim & Audience slice & Worked meaning \\
\midrule
Public-status continuity & release-note slice & one public-status route, one continuity branch, shared witness core \\
Public-status continuity & referee slice & one referee route, one continuity branch, the same shared witness core \\
Compact diff recomputation & auditor slice & one auditor route, one recomputation branch, shared support-only witness core \\
\bottomrule
\end{tabularx}
\caption{Audience-specific witness slices keep per-audience route/branch burden minimal while remaining derivable from one aggregate witness pack and one shared witness core.}
\label{tab:witnessslices}
\end{table}

\section{Worked example pointer}
The maintained worked bundle now ships \path{example_successor_challenge_coverage_witness_slices.json}.\cite{series:workedexample}
It records ten audience-minimal slices for the canned Case-B successor: two slices each for the two public-status coverage entries and the three compact-diff coverage entries.
Each slice names exactly one audience route, the matching route-specific branch set, the same witness-pack key already declared for that coverage entry, and the same witness-core key reused across audiences.
That keeps Synthesis~39 about \emph{what} is covered, Synthesis~42 about the aggregate route/branch proving bundle for a coverage claim, Synthesis~44 about the \emph{audience-invariant proving core}, and this note about the \emph{audience-minimal route/branch delta} for one audience-specific use of that same claim.\cite{series:challengecoverage,series:challengecoveragewitnesspacks,series:challengecoveragewitnesscores}

\paragraph{Why this helps the archive.}
The series can now give later papers a smallest sufficient audience-specific route/branch burden without making them restate or reconstruct the global certification bundle by hand.
That is the right next tightening for brief referee-grade notes.


\paragraph{A minimal public witness-slice card.}
\begin{table}[t]
\centering
\small
\begin{tabularx}{\linewidth}{@{}lX@{}}
\toprule
Field & Minimal public answer \\
\midrule
Release-note continuity slice & \nolinkurl{public_status_sentence_continuity_piece_coverage_release_note_slice} keeps route \nolinkurl{release_note_public_status_sentence}, branch \nolinkurl{release_note_public_status_sentence_continuity_piece}, coverage entry \nolinkurl{public_status_sentence_continuity_piece_coverage}, witness pack \nolinkurl{public_status_sentence_continuity_piece_witness_pack}, and shared witness core \nolinkurl{public_status_sentence_continuity_piece_witness_core} for the release-note use of the continuity claim \\
Auditor recomputation slice & \nolinkurl{compact_numeric_diff_recompute_piece_coverage_auditor_slice} keeps route \nolinkurl{auditor_compact_diff}, branch \nolinkurl{auditor_compact_diff_recompute_piece}, coverage entry \nolinkurl{compact_numeric_diff_recompute_piece_coverage}, witness pack \nolinkurl{compact_numeric_diff_recompute_piece_witness_pack}, and shared witness core \nolinkurl{compact_numeric_diff_recompute_piece_witness_core} for the auditor use of the support-only recomputation claim \\
Selection rule & Stop here only while the live issue is the audience-specific route-and-branch proving delta for one already-certified coverage claim, not the aggregate witness pack, the audience-invariant proving core, or the later one-piece answer bundle \\
Left/right boundary & Reopen Synthesis~39--42 when the real dispute is the coverage entry itself, an imported surface or terminal owner, or the aggregate pack; reopen Synthesis~44 when the shared core is live; widen to Synthesis~45 only when the honest issue has become the one-piece answer bundle rather than the proving delta \\
Paired-terminal boundary & Once the exact audience-specific proving delta is no longer the moving part, later terse papers should usually stop at Synthesis~54 together with Synthesis~74--75 rather than reopen this note merely to restate a stable answer-side handoff or paired cutover judgment \\
\bottomrule
\end{tabularx}
\caption{A minimal public witness-slice card. This is the smallest audience-specific proving object later terse papers can quote directly when the live issue is one maintained route-and-branch delta for one already-certified coverage claim.}
\label{tab:minimal-public-witness-slice-card}
\end{table}

This card is intentionally non-decorative: it pins the actual worked slice keys, the imported route / branch / pack / core anchors they preserve, and the stop rule for one audience-specific proving delta without silently re-owning the aggregate witness pack, the audience-invariant core, or the later answer-side bundle.\cite{series:workedexample}

\paragraph{A forced public witness-slice matrix.}
\begin{table}[t]
\centering
\small
\begin{tabularx}{\linewidth}{@{}lX@{}}
\toprule
Field & Forced public answer \\
\midrule
Release-note continuity floor & \nolinkurl{public_status_sentence_continuity_piece_coverage_release_note_slice} still forces the same audience-specific proving delta only while the same route \nolinkurl{release_note_public_status_sentence}, branch \nolinkurl{release_note_public_status_sentence_continuity_piece}, coverage entry \nolinkurl{public_status_sentence_continuity_piece_coverage}, witness pack \nolinkurl{public_status_sentence_continuity_piece_witness_pack}, and shared witness core \nolinkurl{public_status_sentence_continuity_piece_witness_core} still agree on the release-note use of \nolinkurl{continuity_piece} \\
Auditor recomputation floor & \nolinkurl{compact_numeric_diff_recompute_piece_coverage_auditor_slice} still forces the same audience-specific proving delta only while the same route \nolinkurl{auditor_compact_diff}, branch \nolinkurl{auditor_compact_diff_recompute_piece}, coverage entry \nolinkurl{compact_numeric_diff_recompute_piece_coverage}, witness pack \nolinkurl{compact_numeric_diff_recompute_piece_witness_pack}, and shared witness core \nolinkurl{compact_numeric_diff_recompute_piece_witness_core} still agree on the auditor use of the support-only recomputation piece \\
Pack/core agreement floor & Those shortcuts remain honest only while the slice is still the audience-minimal face of the same aggregate witness pack and the same audience-invariant witness core; a slice does not keep the same citation once its pack or shared core has changed \\
Coverage-mode floor & The same slice key still forces the same proving delta only while its coverage mode remains the same; if a support-only slice becomes surfaced, or a surfaced slice loses the supporting shared core hooks that justified that route, the matrix has already failed closed \\
Staleness trigger floor & This matrix goes stale as soon as the audience changes, the route or route-specific branch set changes, the coverage entry rebinds, the witness pack or shared core changes, or the maintained slice stops being the exact audience-minimal delta for that claim \\
Escalation pointer & Stop here only while one maintained slice still forces the same audience-specific route-and-branch delta; widen to Synthesis~45 when the dispute becomes the one-piece answer bundle, reopen Synthesis~42 when the aggregate pack is live, or reopen Synthesis~44 when the shared proving core itself has drifted \\
\bottomrule
\end{tabularx}
\caption{A forced public witness-slice matrix. This is the smallest audience-specific proving object later terse papers can quote directly when the live issue is whether one maintained slice still forces the same route-and-branch delta rather than merely which slice exists in the abstract.}
\label{tab:forced-public-witness-slice-matrix}
\end{table}

This matrix is intentionally non-decorative: it pins the actual worked slice keys that still force one exact audience-specific proving delta, together with the conditions under which that forcing remains honest or goes stale, without silently re-owning the aggregate pack, the shared proving core, or the later answer-side exact-bundle owners.\cite{series:workedexample}

\paragraph{One tiny witness-slice drill.}
For the worked release-note continuity use, \nolinkurl{public_status_sentence_continuity_piece_coverage_release_note_slice} still forces the same proving delta while the same release-note route, continuity branch, coverage entry, witness pack, and shared witness core still converge on the release-note use of \nolinkurl{continuity_piece}. If the live question widens from ``is this still the release-note proving delta for the continuity claim?'' to ``what is the whole one-piece answer bundle?'' the honest stop is no longer this note but Synthesis~45. If the live question stays audience-specific but the branch set widens, the audience changes, or the maintained slice no longer points to the same aggregate pack and shared core, the matrix has gone stale and a later terse paper must reopen the exact witness-slice owner rather than keep quoting the old delta. Only once the audience-specific proving delta is fixed and the remaining issue becomes the maintained answer-side handoff or paired terminal cutover should the paper leave this note for Synthesis~54 or Synthesis~74--75.\cite{series:workedexample}

\begin{thebibliography}{99}\small

\bibitem{series:challengecoverage}
Anonymous.
\newblock Challenge Coverage Certificates for Successor Maintenance: Audience Coverage and Support-Only Hooks.
\newblock Series synthesis note (Synthesis~39), 2026. (This archive.)

\bibitem{series:challengesurfacehooks}
Anonymous.
\newblock Challenge Surface Hooks for Successor Maintenance: Piece-to-Segment Realization and Shared Surface Bindings.
\newblock Series synthesis note (Synthesis~40), 2026. (This archive.)

\bibitem{series:challengeterminalwitnesses}
Anonymous.
\newblock Challenge Terminal Witness Ledgers for Successor Maintenance: Narrow Stop Ownership for Piece Hooks and Support-Only Escalation.
\newblock Series synthesis note (Synthesis~41), 2026. (This archive.)

\bibitem{series:challengecoveragewitnesspacks}
Anonymous.
\newblock Challenge Coverage Witness Packs for Successor Maintenance: Aggregate Route-and-Branch Certification Dependencies for One Coverage Claim.
\newblock Series synthesis note (Synthesis~42), 2026. (This archive.)

\bibitem{series:challengecoveragewitnesscores}
Anonymous.
\newblock Challenge Coverage Witness Cores for Successor Maintenance: Audience-Invariant Certification Dependencies for One Coverage Claim.
\newblock Series synthesis note (Synthesis~44), 2026. (This archive.)

\bibitem{series:challengeanswerreviewmenus}
Anonymous.
\newblock Challenge Answer Review Menus for Successor Maintenance: Stable Request Classes and Terminal Citation Normal Form.
\newblock Series synthesis note (Synthesis~54), 2026. (This archive.)

\bibitem{series:pairedterminalcitationdiscipline}
Anonymous.
\newblock Paired Terminal Citation Discipline for Review Spines: Stable Joint Reuse of the Checked-Answer Stop and the Request-Side Terminal Stop.
\newblock Series synthesis note (Synthesis~74), 2026. (This archive.)

\bibitem{series:pairedterminalcutoverrules}
Anonymous.
\newblock Paired Terminal Cutover Rules for Review Spines: When Later Terse Papers Stop, Widen, or Reopen.
\newblock Series synthesis note (Synthesis~75), 2026. (This archive.)

\bibitem{series:workedexample}
Anonymous.
\newblock Worked Example for Anonymous-DHT Receipts: Minimal Public Surface, Cross-Release Diff, and Successor Interlock.
\newblock Series synthesis note (Synthesis~17), 2026. (This archive.)

\end{thebibliography}
\end{document}
